% Options for packages loaded elsewhere
\PassOptionsToPackage{unicode}{hyperref}
\PassOptionsToPackage{hyphens}{url}
\PassOptionsToPackage{dvipsnames,svgnames,x11names}{xcolor}
\documentclass[
  10pt,
]{ctexart}
\usepackage{xcolor}
\usepackage[margin=2cm]{geometry}
\usepackage{amsmath,amssymb}
\setcounter{secnumdepth}{-\maxdimen} % remove section numbering
\usepackage{iftex}
\ifPDFTeX
  \usepackage[T1]{fontenc}
  \usepackage[utf8]{inputenc}
  \usepackage{textcomp} % provide euro and other symbols
\else % if luatex or xetex
  \usepackage{unicode-math} % this also loads fontspec
  \defaultfontfeatures{Scale=MatchLowercase}
  \defaultfontfeatures[\rmfamily]{Ligatures=TeX,Scale=1}
\fi
\usepackage{lmodern}
\ifPDFTeX\else
  % xetex/luatex font selection
  \setmainfont[]{DejaVu Sans}
\fi
% Use upquote if available, for straight quotes in verbatim environments
\IfFileExists{upquote.sty}{\usepackage{upquote}}{}
\IfFileExists{microtype.sty}{% use microtype if available
  \usepackage[]{microtype}
  \UseMicrotypeSet[protrusion]{basicmath} % disable protrusion for tt fonts
}{}
\makeatletter
\@ifundefined{KOMAClassName}{% if non-KOMA class
  \IfFileExists{parskip.sty}{%
    \usepackage{parskip}
  }{% else
    \setlength{\parindent}{0pt}
    \setlength{\parskip}{6pt plus 2pt minus 1pt}}
}{% if KOMA class
  \KOMAoptions{parskip=half}}
\makeatother
\setlength{\emergencystretch}{3em} % prevent overfull lines
\providecommand{\tightlist}{%
  \setlength{\itemsep}{0pt}\setlength{\parskip}{0pt}}
\usepackage{bookmark}
\IfFileExists{xurl.sty}{\usepackage{xurl}}{} % add URL line breaks if available
\urlstyle{same}
% fallback for those not using the hyperref driver hyperxmp:
\makeatletter
\@ifundefined{xmpquote}{\newcommand{\xmpquote}[1]{#1}}{}
\makeatother
\hypersetup{
  colorlinks=true,
  linkcolor={Maroon},
  filecolor={Maroon},
  citecolor={Blue},
  urlcolor={Blue},
  pdfcreator={LaTeX via pandoc}}

\author{}
\date{}

\usepackage{pdflscape,fvextra}
\setlength{\tabcolsep}{3pt}
\begin{document}

\section{LUAD Top-10
多智能体证据审阅（v1）}\label{luad-top-10-ux591aux667aux80fdux4f53ux8bc1ux636eux5ba1ux9605v1}

后续对本版 65
条引文所做的\href{evidence_scope_audit.md}{来源与药物指向预筛}将 23
条列为需要进一步核对；该结果不改写这里的冻结分级。

\begin{quote}
\textbf{这是文献分诊（literature triage），不是疗效证据。} 分级只描述在
PubMed
中检索到、并且通过引用校验的文献现状，用来安排后续研究的优先级；不构成治疗建议，也不证明任何候选药物对肺腺癌有效。
\end{quote}

运行日期：2026-09-28。模型：\texttt{deepseek-flash}（DeepSeek beta
端点，strict function calling）。
代码：\texttt{src/\allowbreak{}drug\_\allowbreak{}repurposing\_\allowbreak{}agent/\allowbreak{}multi\_\allowbreak{}agent\_\allowbreak{}review.\allowbreak{}py}；运行器：\texttt{scripts/\allowbreak{}run\_\allowbreak{}multi\_\allowbreak{}agent\_\allowbreak{}review.\allowbreak{}py}。
结果：\texttt{benchmark/\allowbreak{}results/\allowbreak{}multi\_\allowbreak{}agent\_\allowbreak{}review\_\allowbreak{}v1.\allowbreak{}json}。 入口
trace：\texttt{artifacts/\allowbreak{}multi\_\allowbreak{}agent\_\allowbreak{}review/\allowbreak{}traces/\allowbreak{}}；每次模型请求的
trace：\texttt{artifacts/\allowbreak{}multi\_\allowbreak{}agent\_\allowbreak{}review/\allowbreak{}provider\_\allowbreak{}traces/\allowbreak{}}（均被
git 忽略）。 输入：\texttt{configs/\allowbreak{}luad\_\allowbreak{}top10\_\allowbreak{}evidence\_\allowbreak{}v1.\allowbreak{}json}，其
SHA-256 记录在结果中。

\subsection{流程}\label{ux6d41ux7a0b}

\begin{enumerate}
\def\labelenumi{\arabic{enumi}.}
\tightlist
\item
  \textbf{Literature Agent}：沿用现有的 \texttt{PubMedClient}（NCBI
  E-utilities，限速 0.36 s），并扩展了 ESearch ID 检索和 EFetch
  摘要获取。检索式为
  \texttt{"\textless{}药名\textgreater{}"{[}Title/\allowbreak{}Abstract{]}\ AND\ (lung\ adenocarcinoma\ OR\ non-small\ cell\ lung\ cancer\ OR\ NSCLC\ OR\ A549)}，按相关度取前
  8 篇，只保留有摘要的记录。这个 Agent
  提出支持性论断，每条论断都必须对应一个 PMID，并附上摘要中的原句。
\item
  \textbf{Critic Agent}：另做一次检索，在同样的药名和肺癌词之外，再加上
  resistance/toxicity/adverse/"no effect"/"did
  not"/failed/worse/reduced/survival
  等负面词。对于糖皮质激素类候选药（通路中含
  glucocorticoid），还会补充一次共享的类别级检索（糖皮质激素/地塞米松 +
  NSCLC + 耐药/免疫治疗，取前 5 篇）。Critic
  要对每一条已验证的支持论断给出 stands/weakened/refuted
  的判断，并列出反证，反证同样要附原文引用。
\item
  \textbf{确定性校验器}：引用的 PMID 必须出现在该 Agent
  实际检索到的记录集合中；引用句在空白归一化后，必须是该 PMID
  摘要的逐字子串，且长度至少 20
  个字符。不合格的论断会被丢弃并计数。Critic 引用的、不存在的 claim\_id
  也会被丢弃。
\item
  \textbf{Coordinator Agent}：从 \{SUPPORTED,
  PROMISING\_BUT\_INCOMPLETE, CONFLICTING, INSUFFICIENT\_EVIDENCE\}
  中选出一个分级，并写一段简短理由（strict
  JSON）。之后执行确定性护栏，护栏\textbf{只能降级}：没有任何已验证的支持论断时一律为
  INSUFFICIENT\_EVIDENCE；SUPPORTED 至少需要 2
  个未被驳斥、候选药物本身范围内的
  PMID。如果某个候选药物既没有检索到记录，也没有任何论断，就直接判为
  INSUFFICIENT\_EVIDENCE，不调用模型。
\end{enumerate}

\subsection{结果}\label{ux7ed3ux679c}

\textbf{排名 1 · hydrocortisone}

\begin{itemize}
\tightlist
\item
  \textbf{支持检索命中 / 有摘要：} 61 / 8
\item
  \textbf{已验证支持论断：} 6
\item
  \textbf{Critic 判定：} 6 条均为 weakened
\item
  \textbf{已验证反证：} 6
\item
  \textbf{多智能体分级：} INSUFFICIENT\_EVIDENCE
\item
  \textbf{现有单轮结论：} insufficient\_evidence
\end{itemize}

\textbf{排名 2 · beclomethasone-dipropionate}

\begin{itemize}
\tightlist
\item
  \textbf{支持检索命中 / 有摘要：} 7 / 7
\item
  \textbf{已验证支持论断：} 6
\item
  \textbf{Critic 判定：} 5 条 weakened，1 条 stands
\item
  \textbf{已验证反证：} 7
\item
  \textbf{多智能体分级：} INSUFFICIENT\_EVIDENCE
\item
  \textbf{现有单轮结论：} insufficient\_evidence
\end{itemize}

\textbf{排名 3 · clocortolone-pivalate}

\begin{itemize}
\tightlist
\item
  \textbf{支持检索命中 / 有摘要：} 0 / 0
\item
  \textbf{已验证支持论断：} 0
\item
  \textbf{Critic 判定：} ---
\item
  \textbf{已验证反证：} 4（类别级）
\item
  \textbf{多智能体分级：} INSUFFICIENT\_EVIDENCE
\item
  \textbf{现有单轮结论：} insufficient\_evidence
\end{itemize}

\textbf{排名 4 · mometasone}

\begin{itemize}
\tightlist
\item
  \textbf{支持检索命中 / 有摘要：} 6 / 6
\item
  \textbf{已验证支持论断：} 6
\item
  \textbf{Critic 判定：} 6 条均为 weakened
\item
  \textbf{已验证反证：} 4
\item
  \textbf{多智能体分级：} INSUFFICIENT\_EVIDENCE
\item
  \textbf{现有单轮结论：} insufficient\_evidence
\end{itemize}

\textbf{排名 5 · flumetasone}

\begin{itemize}
\tightlist
\item
  \textbf{支持检索命中 / 有摘要：} 0 / 0
\item
  \textbf{已验证支持论断：} 0
\item
  \textbf{Critic 判定：} ---
\item
  \textbf{已验证反证：} 3（类别级）
\item
  \textbf{多智能体分级：} INSUFFICIENT\_EVIDENCE
\item
  \textbf{现有单轮结论：} insufficient\_evidence
\end{itemize}

\textbf{排名 6 · BRD-K84203638}

\begin{itemize}
\tightlist
\item
  \textbf{支持检索命中 / 有摘要：} 0 / 0
\item
  \textbf{已验证支持论断：} 0
\item
  \textbf{Critic 判定：} ---
\item
  \textbf{已验证反证：} 0
\item
  \textbf{多智能体分级：}
  INSUFFICIENT\_EVIDENCE（确定性判定，未调用模型）
\item
  \textbf{现有单轮结论：} insufficient\_evidence
\end{itemize}

\textbf{排名 7 · hydrocortisone-hemisuccinate}

\begin{itemize}
\tightlist
\item
  \textbf{支持检索命中 / 有摘要：} 0 / 0
\item
  \textbf{已验证支持论断：} 0
\item
  \textbf{Critic 判定：} ---
\item
  \textbf{已验证反证：} 4（类别级）
\item
  \textbf{多智能体分级：} INSUFFICIENT\_EVIDENCE
\item
  \textbf{现有单轮结论：} insufficient\_evidence
\end{itemize}

\textbf{排名 8 · fluorometholone}

\begin{itemize}
\tightlist
\item
  \textbf{支持检索命中 / 有摘要：} 0 / 0
\item
  \textbf{已验证支持论断：} 0
\item
  \textbf{Critic 判定：} ---
\item
  \textbf{已验证反证：} 4（类别级）
\item
  \textbf{多智能体分级：} INSUFFICIENT\_EVIDENCE
\item
  \textbf{现有单轮结论：} insufficient\_evidence
\end{itemize}

\textbf{排名 9 · fluticasone}

\begin{itemize}
\tightlist
\item
  \textbf{支持检索命中 / 有摘要：} 24 / 8
\item
  \textbf{已验证支持论断：} 6
\item
  \textbf{Critic 判定：} 4 条 weakened，2 条 stands
\item
  \textbf{已验证反证：} 7
\item
  \textbf{多智能体分级：} INSUFFICIENT\_EVIDENCE
\item
  \textbf{现有单轮结论：} insufficient\_evidence
\end{itemize}

\textbf{排名 10 · diflorasone}

\begin{itemize}
\tightlist
\item
  \textbf{支持检索命中 / 有摘要：} 0 / 0
\item
  \textbf{已验证支持论断：} 0
\item
  \textbf{Critic 判定：} ---
\item
  \textbf{已验证反证：} 2（类别级）
\item
  \textbf{多智能体分级：} INSUFFICIENT\_EVIDENCE
\item
  \textbf{现有单轮结论：} insufficient\_evidence
\end{itemize}

\textbf{与现有单轮结论的比较}：10/10 一致，全部为
INSUFFICIENT\_EVIDENCE。现有单轮结论来自
\texttt{configs/\allowbreak{}luad\_\allowbreak{}top10\_\allowbreak{}evidence\_\allowbreak{}v1.\allowbreak{}json} 和
\texttt{docs/\allowbreak{}luad\_\allowbreak{}top10\_\allowbreak{}evidence\_\allowbreak{}matrix.\allowbreak{}md}。没有任何一个分级是由护栏改出来的：每个调用了模型的候选，Coordinator
给出的都已经是 INSUFFICIENT\_EVIDENCE。

各候选的主要理由（由 Coordinator 生成，已经过校验）：

\begin{itemize}
\tightlist
\item
  \textbf{hydrocortisone}：检索到的支持记录几乎都是免疫检查点抑制剂相关的垂体炎/肾上腺功能不全个案报告，氢化可的松在其中是替代治疗，不涉及抗肿瘤作用。
\item
  \textbf{beclomethasone-dipropionate、mometasone、fluticasone}：支持论断全部来自
  A549 细胞上的体外抗炎、转录抑制或 CYP3A5 诱导等替代指标，没有任何
  LUAD/NSCLC 抗肿瘤终点、动物实验或临床证据。
\item
  \textbf{其余 6
  个}：按确切药名没有检索到候选药本身的文献，只有糖皮质激素类别级的负面信号（免疫治疗干扰、耐药）。
\end{itemize}

\textbf{幻觉引用拒绝率}：两个 Agent 共提出 65 条带引用的论断（支持 24
条，反证 41 条），校验器拒绝 \textbf{0 条（0.0\%）}。在 strict
工具调用、并且把检索到的摘要全文提供给模型的条件下，DeepSeek 没有编造
PMID，也没有改写引文。

调用量：共 22 次模型调用（4 个有支持文献的候选各 3 次
Literature+Critic+Coordinator，5 个只有类别级记录的候选各 2 次
Critic+Coordinator，BRD-K84203638 为 0 次），全部成功，0 次重试。共约
6.3 万 token，估算费用为 0.014 美元（off-peak）至 0.028
美元（peak）。PubMed 请求另计，不收费。

\subsection{局限与观察}\label{ux5c40ux9650ux4e0eux89c2ux5bdf}

\begin{itemize}
\tightlist
\item
  \textbf{逐字引用不等于引用正确}。校验器只能保证引文确实存在于该 PMID
  的摘要中，无法保证引文能支持对应的论断。人工抽查发现：Critic
  曾把一篇关于柚皮素纳米载体的 A549 摘要（PMID 33249629）标为
  fluticasone 的 \texttt{candidate} 级反证；多次把一条勘误声明（PMID
  41376996，Published Erratum，摘要只有一句``This corrects the
  article''）当作证据引用；还常用"X
  是安全/有益的"这种稻草人式表述来写反证。所以反证数量不能当作真实的矛盾强度，下一步应该加入论断--引文蕴含（NLI）检查或人工复核。
\item
  \textbf{检索式与既有矩阵不同}。hydrocortisone
  在既有矩阵中的候选药本身文献（PMID
  1533045、35676421）没有进入按相关度排序的前 8 篇；fluticasone
  的预防性荟萃分析（PMID 35843928）也没有被检索到。beclomethasone
  的两篇既有文献（12538830、24555085）都检索到了。这是一次有边界的检索快照，不能证明"不存在证据"。
\item
  只按确切药名检索，没有扩展同义词、盐或酯形式；身份问题（立体化学、剂型歧义）只作为
  Coordinator 的输入，没有经过化学层面的核验。
\item
  只用了一个模型，每个候选只跑一轮，结果没有做重复性检验。
\end{itemize}

\subsection{复现}\label{ux590dux73b0}

\begin{Verbatim}[breaklines=true,breakanywhere=true]
python scripts/run_multi_agent_review.py   # 需要 .env 中的 DEEPSEEK_API_KEY；已完成的候选会从 artifacts/multi_agent_review/checkpoint_v1.jsonl 续跑
\end{Verbatim}

\end{document}
